\chapter{Analiza problemu i przegląd rozwiązań}\label{ch:analiza}
% Ok. 8-10 stron. Wymaga zebrania źródeł i zdjęć z rozbiórki.

\section{Konstrukcja IBM Model M SSK}\label{sec:modelm}
% - Buckling spring over membrane, stalowa płyta, układ SSK, oryginalny kontroler.
% - Membrana: dwie warstwy przewodzące + spacer; ogony; bolt-mod. Zdjęcia własne.
% - Źródło: patent IBM US 4,118,611 (buckling spring).
\todo{Napisać.}

\section{Matryca membranowa: ghosting, masking, NKRO}\label{sec:matryca}
% - Skanowanie wiersz/kolumna: 6 + 17 = 23 GPIO zamiast 88 wejść.
% - Rysunek prostokąta 2x2: ghosting, masking; dioda na klawisz daje NKRO.
% - Oryginalna membrana bez diod -> decyzja o własnej membranie.
\todo{Napisać.}

\section{Istniejące kontrolery zastępcze i ich porównanie}\label{sec:kontrolery}
% - Oryginalny kontroler, otwarte zamienniki (Teensy/ATmega32U4), FSSK (reference/fssk-v1.00).
% - Tabela: MCU, przewód/bezprzewodowo, NKRO, ingerencja w obudowę, otwartość.
% - Wniosek: brak bezprzewodowego drop-in z NKRO.
\todo{Napisać.}

\section{Łączność bezprzewodowa w klawiaturach}\label{sec:ble}
% - BLE HID over GATT, interwał połączenia vs opóźnienie vs prąd; USB HID; dongle 2,4 GHz krótko.
\todo{Napisać.}

\section{Oprogramowanie klawiatur: QMK, ZMK}\label{sec:qmkzmk}
% - QMK vs ZMK (Zephyr, BLE, devicetree, MIT); uzasadnienie wyboru: HD 'Firmware: ZMK na Zephyr 4.1'.
\todo{Napisać.}

